% ============================================================
%  GUÍA 3: POTENCIACIÓN DE NÚMEROS ENTEROS
%  Curso Introductorio de Matemáticas — 1er Año
%  Autor: Ing. Gabriel Astudillo
%  Sesiones cubiertas: Semana 3 (clases 9 a 12)
% ============================================================

\documentclass[12pt, letterpaper]{article}

% ── Paquetes ─────────────────────────────────────────────────

\usepackage{mathwizards-guia}
\mwguia{Guía 3 — Potenciación de Números Enteros}{Ing. Gabriel Astudillo $\cdot$ Curso 1er Año}

% ── Comandos propios de esta guía ───────────────────
\newcommand{\Z}{\mathbb{Z}}

\begin{document}
% ============================================================

% ── Portada / Título ─────────────────────────────────────────
\mwportada{Guía de Estudio 3}{Potenciación de Números Enteros}{Definición, Propiedades y Primeros Pasos en Álgebra}

\vspace{4pt}
\begin{cuadroinfo}
  \small
  \textbf{Presentación:} \textquestiondown Alguna vez has escuchado que «una hoja de papel se puede doblar a la mitad tantas veces que su espesor crece rapidísimo»? Eso se explica con las \textbf{potencias}: una forma súper compacta de escribir multiplicaciones repetidas. Si doblas una hoja $5$ veces, aparecen $2^5 = 32$ capas; \textquestiondown increíble, verdad? En esta guía aprenderás qué es una potencia, sus propiedades secretas y, al final, las usaremos con \emph{letras} (eso es el comienzo del \textbf{álgebra}).
  \\[4pt]
  \textbf{Nivel de Dificultad:}\quad
  \facil{} Básico / Directo \quad
  \medio{} Intermedio \quad
  \dificil{} Desafiante
\end{cuadroinfo}

\vspace{8pt}

% ============================================================
\section{\textquestiondown Qué es una Potencia?}
% ============================================================

\subsection{Base, exponente y potencia}

Multiplicar un número por sí mismo varias veces es muy común, así que los matemáticos inventaron una forma corta de escribirlo.

\begin{cuadroteoria}
  \textbf{Definición:} una \textbf{potencia} es una multiplicación repetida. Se escribe $a^n$ y se lee «$a$ elevado a la $n$».
  \[
    a^n = \underbrace{a \times a \times a \times \cdots \times a}_{n \text{ veces}}
  \]
  \begin{itemize}
    \item \textbf{Base} ($a$): el número que se multiplica.
    \item \textbf{Exponente} ($n$): cuántas veces se multiplica.
    \item \textbf{Potencia}: el resultado.
  \end{itemize}
\end{cuadroteoria}

\noindent\textbf{Ejemplos:}
\begin{itemize}[leftmargin=*]
  \item $2^3 = 2 \times 2 \times 2 = 8$ (se lee «dos elevado a la tres» o «dos al cubo»).
  \item $3^2 = 3 \times 3 = 9$ (se lee «tres elevado al cuadrado» o «tres al cuadrado»).
  \item $10^3 = 10 \times 10 \times 10 = 1000$ (mil).
\end{itemize}

\newpage

\begin{cuadroadvertencia}
  \textbf{Exponentes especiales:}
  \begin{itemize}
    \item Cualquier número elevado a la $1$ es él mismo: $a^1 = a$. Ejemplo: $7^1 = 7$.
    \item Cualquier número distinto de cero elevado a la $0$ es $1$: $a^0 = 1$. Ejemplo: $5^0 = 1$. (\textquestiondown Por qué? Porque al dividir $5^1 \div 5^1 = 1$ y también $5^{1-1} = 5^0$... ya lo verás en las propiedades.)
  \end{itemize}
\end{cuadroadvertencia}

\subsection{El signo del resultado}

\begin{cuadroteoria}
  \textbf{Regla del signo de una potencia:}
  \begin{itemize}
    \item Base \textbf{positiva} $\Rightarrow$ resultado siempre \textbf{positivo}. ($3^2 = 9$, $2^3 = 8$.)
    \item Base \textbf{negativa} con exponente \textbf{par} $\Rightarrow$ resultado \textbf{positivo}: $(-2)^2 = 4$.
    \item Base \textbf{negativa} con exponente \textbf{impar} $\Rightarrow$ resultado \textbf{negativo}: $(-2)^3 = -8$.
  \end{itemize}
\end{cuadroteoria}

\begin{cuadroadvertencia}
  \textbf{¡Cuidado con el paréntesis!} No es lo mismo $(-2)^2$ que $-2^2$:
  \begin{itemize}
    \item $(-2)^2 = (-2) \times (-2) = 4$ (el paréntesis incluye el signo).
    \item $-2^2 = -(2 \times 2) = -4$ (solo el $2$ está al cuadrado, el signo menos queda afuera).
  \end{itemize}
\end{cuadroadvertencia}

\noindent\textbf{Ejemplo resuelto:} \textquestiondown Cuánto es $(-3)^3$?

Como el exponente es impar, el resultado es negativo: $(-3)^3 = (-3) \times (-3) \times (-3) = 9 \times (-3) = -27$.

\subsection{Ejercicios: \textquestiondown Qué es una Potencia?}

\begin{enumerate}[leftmargin=*]
  \item \facil{} $3^2$

  \item \facil{} $2^3$

  \item \facil{} $5^2$

  \item \facil{} $10^3$

  \item \facil{} $1^5$

  \item \facil{} $4^2$

  \item \medio{} $(-2)^2$

  \item \medio{} $(-3)^2$

  \item \medio{} $(-2)^3$

  \item \medio{} $2^4$

  \item \medio{} $7^2$

  \item \medio{} $3^3$

  \item \medio{} $-2^2$ (sin paréntesis). \textquestiondown En qué se diferencia de $(-2)^2$?

  \item \medio{} $(-1)^{201}$ (el exponente es impar... \textquestiondown qué crees que da?)

  \item \medio{} $(-5)^2$

  \item \dificil{} Encuentra la base: $n^2 = 49$. \textquestiondown Cuáles son las dos posibilidades en $\Z$?

  \item \dificil{} Encuentra el exponente: $2^n = 16$.

  \item \dificil{} Una hoja de papel se dobla a la mitad una y otra vez. Cada doblez duplica las capas. \textquestiondown Cuántas capas hay después de doblarla $5$ veces?

  \item \dificil{} Un cultivo de bacterias se duplica cada hora. Si empezamos con una bacteria, \textquestiondown cuántas habrá después de $6$ horas?

  \item \dificil{} \textbf{:} Calcula $(-1)^6 + (-1)^7 + (-1)^8$. (Pista: mira el exponente de cada uno antes de sumar.)
\end{enumerate}

%\newpage

% ============================================================
\section{Propiedades de las Potencias}
% ============================================================

\subsection{Las cinco propiedades mágicas}

Con estas propiedades puedes \emph{comprimir} y \emph{simplificar} potencias sin calcular todo:

\begin{cuadroteoria}
  \textbf{Propiedades de las potencias (con $a$, $b$ enteros y $m$, $n$ naturales):}
  \begin{enumerate}
    \item \textbf{Producto de potencias de igual base:} $a^m \times a^n = a^{m+n}$. \quad Ejemplo: $2^3 \times 2^4 = 2^{3+4} = 2^7$.
    \item \textbf{Cociente de potencias de igual base:} $a^m \div a^n = a^{m-n}$ (con $a \neq 0$). \quad Ejemplo: $3^5 \div 3^2 = 3^{5-2} = 3^3$.
    \item \textbf{Potencia de una potencia:} $(a^m)^n = a^{m \times n}$. \quad Ejemplo: $(2^3)^2 = 2^{3 \times 2} = 2^6$.
    \item \textbf{Potencia de un producto:} $(a \times b)^n = a^n \times b^n$. \quad Ejemplo: $(2 \times 5)^2 = 2^2 \times 5^2 = 4 \times 25 = 100$.
    \item \textbf{Potencia de un cociente:} $(a \div b)^n = a^n \div b^n$ (con $b \neq 0$). \quad Ejemplo: $(10 \div 5)^3 = 10^3 \div 5^3 = 1000 \div 125 = 8$.
  \end{enumerate}
\end{cuadroteoria}

\noindent\textbf{¿Por qué funciona?} Mira el producto: $2^3 \times 2^2 = (2 \times 2 \times 2) \times (2 \times 2)$. \textquestiondown Cuántos doses hay? ¡$5$! Por eso $2^{3+2}$.

\begin{cuadroadvertencia}
  \textbf{Recordatorio:} la propiedad 1 dice que las bases deben ser \emph{iguales}. No puedes usar $2^3 \times 5^2$ como si fueran iguales: ahí se aplica la propiedad 4 (potencia de un producto) solo si \emph{también} los exponentes son iguales.
\end{cuadroadvertencia}

\noindent\textbf{Ejemplo resuelto:} Simplifica $\dfrac{2^5 \times 2^3}{2^4}$.

Arriba: $2^5 \times 2^3 = 2^{5+3} = 2^8$. Luego $2^8 \div 2^4 = 2^{8-4} = 2^4 = 16$.

\subsection{Ejercicios: Propiedades de las Potencias}

\begin{enumerate}[leftmargin=*]
  \item \facil{} $2^2 \times 2^3$

  \item \facil{} $3^3 \times 3^2$

  \item \facil{} $5^4 \times 5^1$

  \item \facil{} $2^5 \div 2^2$

  \item \facil{} $7^3 \div 7^3$

  \item \facil{} $(2^2)^3$

  \item \medio{} $10^3 \times 10^2$

  \item \medio{} $10^6 \div 10^2$

  \item \medio{} $(3^2)^2$

  \item \medio{} $(2 \times 5)^2$

  \item \medio{} $2^3 \times 2^4 \div 2^2$

  \item \medio{} $(5^2)^3 \div 5^4$

  \item \medio{} $6^5 \div 6^2 \times 6^3$ (escríbelo como una sola potencia)

  \item \medio{} $(3 \times 4)^2$

  \item \medio{} $10^5 \div 10^3 \times 10^2$

  \item \dificil{} Simplifica: $(2^3 \times 2^2) \div (2^2 \times 2)$ y calcula el resultado.

  \item \dificil{} Encuentra el exponente: $3^n \times 3^2 = 3^7$.

  \item \dificil{} $8^2 \div 2^2$. (Pista: usa la propiedad del cociente: $8 \div 2 = 4$.)

  \item \dificil{} Un cubo tiene $10$ cm de lado. \textquestiondown Cuál es su volumen? (Recuerda: volumen $=$ lado$^3$.) \textquestiondown Y cuál es el área de una de sus caras (lado$^2$)?

  \item \dificil{} \textbf{:}
        \begin{enumerate}[label=\alph*)]
          \item Calcula $2^{10} \times 5^{10}$ (pista: multiplica las bases).
          \item Simplifica: $\big((2^3)^2 \times 2\big) \div 2^5$ y calcula el resultado.
        \end{enumerate}
\end{enumerate}

\newpage

% ============================================================
\section{Potencias con Letras (Mini Álgebra)}
% ============================================================

\subsection{La base puede ser una letra}

En matemáticas usamos \textbf{letras} para representar números que no conocemos (o que pueden cambiar). Una letra se comporta exactamente como un número.

\begin{cuadroteoria}
  \textbf{Definición:} si $a$ es un número cualquiera, entonces
  \[
    a^3 = a \times a \times a, \qquad a^2 = a \times a, \qquad a^1 = a, \qquad a^0 = 1 \ (a \neq 0).
  \]
  Para \textbf{calcular su valor} solo hay que reemplazar la letra por el número que nos digan.
\end{cuadroteoria}

\noindent\textbf{Ejemplo resuelto 1:} Calcula $a^3$ si $a = 2$.

\[
  a^3 = 2^3 = 2 \times 2 \times 2 = 8
\]

\noindent\textbf{Ejemplo resuelto 2:} Calcula $3a^2$ si $a = 2$.

\textquestiondown Ojo! Primero la potencia, después la multiplicación: $3a^2 = 3 \times a^2 = 3 \times 2^2 = 3 \times 4 = 12$. (Si lo hubiéramos hecho mal, $3 \times 2 = 6$ y $6^2 = 36$... \textquestiondown muy diferente, verdad? Por eso el orden importa.)

\begin{cuadroadvertencia}
  \textbf{No confundir:}
  \begin{itemize}
    \item $2a$ significa $2 \times a$ (multiplicar por $2$).
    \item $a^2$ significa $a \times a$ (multiplicar la letra por sí misma).
    \item $2a \neq a^2$ en general. Ejemplo con $a = 5$: $2a = 10$ pero $a^2 = 25$.
  \end{itemize}
\end{cuadroadvertencia}

\noindent\textbf{Ejemplo resuelto 3:} \textquestiondown Cuál es el área de un cuadrado de lado $x$? El área es $x^2$. Si $x = 7$ cm, el área es $7^2 = 49$ cm$^2$.

% ============================================================
\subsection{Productos y potencias con letras}
% ============================================================

Las propiedades de las potencias que aprendiste en la sección anterior funcionan \emph{igual} cuando la base o los exponentes son letras. Pero hay un truco nuevo: los exponentes también pueden ser \textbf{expresiones con letras}, como $7x$ o $5m$.

\begin{cuadroteoria}
  \textbf{Reglas de oro (con $a \neq 0$ y exponentes que pueden ser letras):}
  \begin{itemize}
    \item \textbf{Producto de igual base:} $a^m \cdot a^n = a^{m+n}$. Los exponentes se \emph{suman}.
          \quad Ejemplo: $a^6 \cdot a^3 = a^{6+3} = a^9$.
    \item \textbf{Potencia de una potencia:} $(a^m)^n = a^{m \times n}$. Los exponentes se \emph{multiplican}.
          \quad Ejemplo: $(p^5)^6 = p^{5 \times 6} = p^{30}$.
  \end{itemize}
  \small\textbf{Importante:} si el exponente es $7x$, lo tratamos como un \emph{único} número: $a^{7x} \cdot a^{6x} = a^{7x + 6x} = a^{13x}$.
\end{cuadroteoria}

\begin{cuadroadvertencia}
  \textbf{Ojo con el orden de las operaciones:} primero se resuelven las potencias y \emph{después} los productos.
  \begin{itemize}
    \item $(2x)^3 = 2^3 \cdot x^3 = 8x^3$ (el $3$ se aplica a \emph{todo}: al $2$ y a la $x$).
    \item $3x^0 = 3 \cdot 1 = 3$ (cualquier número elevado a la $0$ es $1$, y queda el $3$).
    \item $(-2p^3)^2 = (-2)^2 \cdot (p^3)^2 = 4 \cdot p^{3 \times 2} = 4p^6$ (exponente par $\Rightarrow$ positivo).
  \end{itemize}
\end{cuadroadvertencia}

\noindent\textbf{Ejemplo resuelto 4:} Simplifica $(3mn^2)^4$.

Aplicamos la potencia a cada factor: $3^4 \cdot m^4 \cdot (n^2)^4 = 81 \cdot m^4 \cdot n^{2 \times 4} = 81m^4n^8$.

\noindent\textbf{Ejemplo resuelto 5:} Simplifica $\big[(3x)^2 \cdot (2x^3)^2\big]^3$.

Primero cada paréntesis: $(3x)^2 = 9x^2$ y $(2x^3)^2 = 4x^{3 \times 2} = 4x^6$. Luego $(9x^2) \cdot (4x^6) = 36x^8$, y al cubo: $(36x^8)^3 = 36^3 x^{8 \times 3} = 46656\, x^{24}$.

% ============================================================
\subsection{Exponentes negativos con letras}
% ============================================================

Un exponente negativo es la forma elegante de escribir una \emph{fracción}. Significa «uno dividido entre la potencia»:

\begin{cuadroteoria}
  \textbf{Definición:} si $a \neq 0$, entonces
  \[
    a^{-n} = \frac{1}{a^n} \qquad \text{y en general} \qquad a^{-n} \neq \text{un número negativo}.
  \]
  Un exponente negativo \textbf{invierte la base} (la pasa al denominador) y luego se potencia. Ejemplos:
  \begin{itemize}
    \item $a^{-5} = \dfrac{1}{a^5}$.
    \item $a^{-5} \cdot a = a^{-5+1} = a^{-4} = \dfrac{1}{a^4}$.
  \end{itemize}
\end{cuadroteoria}

\begin{cuadroadvertencia}
  \textbf{No confundir exponente negativo con base negativa:}
  \begin{itemize}
    \item $(-2)^2 = 4$ (base negativa, exponente par).
    \item $2^{-2} = \dfrac{1}{2^2} = \dfrac{1}{4}$ (base positiva, exponente negativo \textrightarrow recíproco).
    \item $(b^{-2})^{-8} = b^{(-2) \times (-8)} = b^{16}$ (negativo por negativo da positivo).
  \end{itemize}
\end{cuadroadvertencia}

\noindent\textbf{Ejemplo resuelto 6:} Simplifica $\left(\dfrac{b^{-2}}{b^{-3}}\right)^{2}$.

Cociente de igual base: los exponentes se restan: $b^{-2 - (-3)} = b^{-2+3} = b^1 = b$. Luego $(b)^2 = b^2$.

% ============================================================
\subsection{Potencia de un cociente con letras}
% ============================================================

Cuando una fracción completa se eleva a una potencia, el exponente se aplica tanto al numerador como al denominador:

\begin{cuadroteoria}
  \textbf{Potencia de un cociente:} si $b \neq 0$, entonces
  \[
    \left(\frac{a}{b}\right)^n = \frac{a^n}{b^n}
  \]
  Y usando la regla del cociente de igual base: $\dfrac{a^m}{a^n} = a^{m-n}$. Ejemplos:
  \begin{itemize}
    \item $\left(\dfrac{a^{2x}}{a^3}\right)^3 = \dfrac{a^{2x \times 3}}{a^{3 \times 3}} = \dfrac{a^{6x}}{a^9} = a^{6x - 9}$.
    \item $\dfrac{y^2 \cdot (3y^2)^2}{9y^4}$: arriba $(3y^2)^2 = 9y^4$, así que $\dfrac{y^2 \cdot 9y^4}{9y^4} = y^2$ (¡se simplifica todo!).
  \end{itemize}
\end{cuadroteoria}

\begin{cuadroadvertencia}
  \textbf{Cuidado con los exponentes negativos adentro de un cociente:}
  \[
    \left(\frac{a^{2x}}{a^3}\right)^3 = a^{(2x-3) \cdot 3} = a^{6x-9}
  \]
  Primero se simplifica la fracción (restamos exponentes) y \emph{después} se aplica la potencia. No multipliques $2x \times 3$ antes de restar $3$.
\end{cuadroadvertencia}

% ============================================================
\subsection{Cambio de base}
% ============================================================

A veces los números son potencias \emph{escondidas} de una misma base. Para simplificar, conviene \textbf{escribirlos todos en la misma base}:

\begin{cuadroteoria}
  \textbf{Potencias clave para reconocer de memoria:}
  \begin{center}
    \begin{tabular}{ll}
      \toprule
      \textbf{Base} & \textbf{Potencias} \\
      \midrule
      Base $2$ & $4 = 2^2$, \quad $8 = 2^3$, \quad $16 = 2^4$, \quad $32 = 2^5$, \quad $64 = 2^6$, \quad $128 = 2^7$ \\
      Base $3$ & $9 = 3^2$, \quad $27 = 3^3$, \quad $81 = 3^4$ \\
      \bottomrule
    \end{tabular}
  \end{center}
  \textbf{Pasos:} (1) escribe cada base como potencia de la base común, (2) aplica potencia de una potencia, (3) usa las reglas del producto/cociente.
\end{cuadroteoria}

\begin{cuadroadvertencia}
  \textbf{Nota:} esto es \emph{cambio de base} para \emph{simplificar potencias}, no el cambio de base de los \emph{logaritmos} (que verás más adelante). Aquí solo reescribimos $64$ como $2^6$, $128$ como $2^7$, $27$ como $3^3$ y $9$ como $3^2$.
\end{cuadroadvertencia}

\noindent\textbf{Ejemplo resuelto 7:} Simplifica $(27^p \cdot 9^{3p})^6$.

Escribimos todo en base $3$: $27 = 3^3$ y $9 = 3^2$. Entonces $27^p = (3^3)^p = 3^{3p}$ y $9^{3p} = (3^2)^{3p} = 3^{6p}$. El producto: $3^{3p} \cdot 3^{6p} = 3^{9p}$. Al cubo con el $6$: $(3^{9p})^6 = 3^{9p \times 6} = 3^{54p}$.

\subsection{Ejercicios: Potencias con Letras}

\subsubsection*{Bloque A: Sustituir la letra por un número}
\textit{Reemplaza la letra por el valor indicado y calcula. Recuerda el orden de las operaciones.}

\begin{enumerate}[leftmargin=*]
  \item \facil{} $a^2$ si $a = 4$

  \item \facil{} $b^3$ si $b = 2$

  \item \facil{} $x^2$ si $x = 10$

  \item \facil{} $y^3$ si $y = 1$

  \item \facil{} $2a$ si $a = 5$ \quad (y compáralo con $a^2$)

  \item \facil{} $a^0$ si $a = 7$

  \item \medio{} $3a^2$ si $a = 2$

  \item \medio{} $2b^3$ si $b = 3$

  \item \medio{} $x^2 + x$ si $x = 5$

  \item \medio{} $a^3 - a$ si $a = 3$

  \item \medio{} $10a^2$ si $a = 2$

  \item \medio{} $2x \times x$ si $x = 3$ (pista: $2x \times x = 2x^2$)

  \item \medio{} $y^2 + y^2$ si $y = 4$ (pista: $y^2 + y^2 = 2y^2$)

  \item \medio{} $a^2 \times a$ si $a = 3$ (pista: usa la propiedad del producto)

  \item \medio{} $(2a)^2$ si $a = 3$. \textquestiondown Es igual a $2a^2$? Compruébalo.

  \item \dificil{} $x^3 + x^2 + x$ si $x = 2$

  \item \dificil{} $a^2 b$ si $a = 3$ y $b = 2$

  \item \dificil{} $x^2 y^3$ si $x = 2$ y $y = 3$

  \item \dificil{} Un cuadrado tiene lado $x = 7$ cm. \textquestiondown Cuál es su área ($x^2$)?

  \item \dificil{} \textbf{:}
        \begin{enumerate}[label=\alph*)]
          \item Calcula $3a^2 - 2a + 1$ si $a = 4$.
          \item Con $a = 3$ y $b = 2$, calcula $a^2 + b^2$ y también $(a + b)^2$. \textquestiondown Son iguales? \textquestiondown Qué conclusión sacas?
        \end{enumerate}
\end{enumerate}

\subsubsection*{Bloque B: Productos y potencias con letras}
\textit{Simplifica aplicando las propiedades. Los exponentes pueden ser letras o expresiones (como $7x$).}

\begin{enumerate}[leftmargin=*, start=21]
  \item \facil{} $a^6 \cdot a^3 =$

  \item \facil{} $2^3 \cdot 2^2 =$

  \item \facil{} $(p^5)^6 =$

  \item \facil{} $b \cdot b^x =$

  \item \facil{} $(3x)^2 =$

  \item \medio{} $(2x)^3 \cdot 3x^0 =$

  \item \medio{} $(3x)^3 \cdot (3x)^0 =$

  \item \medio{} $(-2p^3)^2 =$

  \item \medio{} $a^{-5} \cdot a =$

  \item \medio{} $a^{7x} \cdot a^{6x} =$

  \item \medio{} $(b^{-2})^{-8} =$

  \item \medio{} $(3mn^2)^4 =$

  \item \dificil{} $(m^{3a} \cdot m^{7a})^3 =$

  \item \dificil{} $\big[(3x)^2 \cdot (2x^3)^2\big]^3 =$
\end{enumerate}

\subsubsection*{Bloque C: Cocientes de potencias con letras}
\textit{Primero simplifica la fracción (resta los exponentes) y después aplica la potencia exterior si la hay.}

\begin{enumerate}[leftmargin=*, start=35]
  \item \medio{} $\left[y^2 \cdot (3y^2)^2\right] \div 9y^4 =$

  \item \medio{} $\left(\dfrac{a^{2x}}{a^3}\right)^3 =$

  \item \medio{} $\left(\dfrac{k^{7t}}{k^{3t}}\right)^{10} =$

  \item \dificil{} $\left(\dfrac{w^{5m}}{w^m}\right)^{-1} =$

  \item \dificil{} $\left(\dfrac{p^{10x}}{p^{6x}}\right)^{-3} =$

  \item \dificil{} $\left(\dfrac{a^{8m} \cdot a^{2m}}{a^{4m}}\right)^6 =$

  \item \dificil{} $\left(\dfrac{x^{4a} \cdot x^{5b}}{x^{2a} \cdot x^{3b}}\right)^6 =$

  \item \dificil{} $\left(\dfrac{n^{5x}}{n^{3x}} \cdot \dfrac{n^{7x}}{n^{3x}}\right)^6 =$
\end{enumerate}

\subsubsection*{Bloque D: Cambio de base y Retos de los Campeones}
\textit{Escribe cada base como potencia de la base común ($64 = 2^6$, $128 = 2^7$, $27 = 3^3$, $9 = 3^2$) y simplifica.}

\begin{enumerate}[leftmargin=*, start=43]
  \item \muyDificil{} $(64^{5x} \div 128^x)^5 =$

  \item \muyDificil{} \textbf{:} $(27^p \cdot 9^{3p})^6 =$
\end{enumerate}

\newpage

% ============================================================
\section{Clave de Respuestas}
% ============================================================

\subsection*{\textquestiondown Qué es una Potencia? (Ejercicios 1--20)}

\begin{enumerate}[leftmargin=*, label=\textbf{\arabic*.}]
  \item $9$
  \item $8$
  \item $25$
  \item $1000$
  \item $1$
  \item $16$
  \item $(-2) \times (-2) = 4$
  \item $(-3) \times (-3) = 9$
  \item $(-2) \times (-2) \times (-2) = -8$
  \item $16$
  \item $49$
  \item $27$
  \item $-4$. Diferencia: en $-2^2$ solo el $2$ está al cuadrado y queda $-(4) = -4$; en $(-2)^2$ el signo está dentro del paréntesis y da $+4$.
  \item $-1$ (exponente impar: negativo)
  \item $25$
  \item $n = 7$ o $n = -7$ (ambos cumplen: $7^2 = 49$ y $(-7)^2 = 49$)
  \item $n = 4$ (porque $2^4 = 16$)
  \item $2^5 = 32$ capas
  \item $2^6 = 64$ bacterias
  \item $(-1)^6 = 1$, $(-1)^7 = -1$, $(-1)^8 = 1$; total: $1 - 1 + 1 = 1$
\end{enumerate}

\subsection*{Propiedades de las Potencias (Ejercicios 21--40)}

\begin{enumerate}[leftmargin=*, label=\textbf{\arabic*.}]
  \setcounter{enumi}{20}
  \item $2^5 = 32$
  \item $3^5 = 243$
  \item $5^5 = 3125$
  \item $2^3 = 8$
  \item $7^0 = 1$
  \item $(2^2)^3 = 2^{2 \times 3} = 2^6 = 64$
  \item $10^5 = 100\,000$
  \item $10^4 = 10\,000$
  \item $(3^2)^2 = 3^4 = 81$
  \item $(2 \times 5)^2 = 10^2 = 100$
  \item $2^3 \times 2^4 = 2^7$; luego $2^7 \div 2^2 = 2^5 = 32$
  \item $(5^2)^3 = 5^6$; luego $5^6 \div 5^4 = 5^2 = 25$
  \item $6^5 \div 6^2 = 6^3$; luego $6^3 \times 6^3 = 6^6$
  \item $(3 \times 4)^2 = 12^2 = 144$
  \item $10^5 \div 10^3 = 10^2$; luego $10^2 \times 10^2 = 10^4 = 10\,000$
  \item $(2^3 \times 2^2) = 2^5$ y $(2^2 \times 2) = 2^3$; luego $2^5 \div 2^3 = 2^2 = 4$
  \item $3^n \times 3^2 = 3^{n+2} = 3^7$, así que $n + 2 = 7$ y $n = 5$
  \item $8^2 \div 2^2 = (8 \div 2)^2 = 4^2 = 16$
  \item Volumen: $10^3 = 1000$ cm$^3$. Área de una cara: $10^2 = 100$ cm$^2$.
  \item a) $2^{10} \times 5^{10} = (2 \times 5)^{10} = 10^{10} = 10\,000\,000\,000$. \quad b) $(2^3)^2 = 2^6$; $2^6 \times 2 = 2^7$; $2^7 \div 2^5 = 2^2 = 4$.
\end{enumerate}

\subsection*{Potencias con Letras (Ejercicios 41--84)}

\begin{enumerate}[leftmargin=*, label=\textbf{\arabic*.}]
  \setcounter{enumi}{40}
  \item $4^2 = 16$
  \item $2^3 = 8$
  \item $10^2 = 100$
  \item $1^3 = 1$
  \item $2a = 2 \times 5 = 10$; \quad $a^2 = 25$ (no son iguales)
  \item $7^0 = 1$
  \item $3 \times 2^2 = 3 \times 4 = 12$
  \item $2 \times 3^3 = 2 \times 27 = 54$
  \item $5^2 + 5 = 25 + 5 = 30$
  \item $3^3 - 3 = 27 - 3 = 24$
  \item $10 \times 2^2 = 10 \times 4 = 40$
  \item $2x^2 = 2 \times 3^2 = 2 \times 9 = 18$
  \item $2y^2 = 2 \times 4^2 = 2 \times 16 = 32$
  \item $a^2 \times a = a^3 = 3^3 = 27$
  \item $(2a)^2 = 6^2 = 36$; \quad $2a^2 = 2 \times 3^2 = 18$. \textquestiondown No son iguales! El paréntesis cambia todo.
  \item $2^3 + 2^2 + 2 = 8 + 4 + 2 = 14$
  \item $a^2 b = 3^2 \times 2 = 9 \times 2 = 18$
  \item $x^2 y^3 = 2^2 \times 3^3 = 4 \times 27 = 108$
  \item $x^2 = 7^2 = 49$ cm$^2$
  \item a) $3(4^2) - 2(4) + 1 = 48 - 8 + 1 = 41$. \quad b) $a^2 + b^2 = 9 + 4 = 13$; $(a + b)^2 = 5^2 = 25$. \textquestiondown No son iguales! Conclusión: $(a + b)^2 \neq a^2 + b^2$ (ese es un error muy común que ya nunca cometerás).
  \item $a^6 \cdot a^3 = a^9$
  \item $2^3 \cdot 2^2 = 2^5 = 32$
  \item $(p^5)^6 = p^{30}$
  \item $b \cdot b^x = b^{1+x}$
  \item $(3x)^2 = 3^2 \cdot x^2 = 9x^2$
  \item $(2x)^3 \cdot 3x^0 = 8x^3 \cdot 3 = 24x^3$
  \item $(3x)^3 \cdot (3x)^0 = 27x^3 \cdot 1 = 27x^3$
  \item $(-2p^3)^2 = (-2)^2 \cdot (p^3)^2 = 4p^6$
  \item $a^{-5} \cdot a = a^{-4} = \dfrac{1}{a^4}$
  \item $a^{7x} \cdot a^{6x} = a^{13x}$
  \item $(b^{-2})^{-8} = b^{(-2)(-8)} = b^{16}$
  \item $(3mn^2)^4 = 3^4 \cdot m^4 \cdot (n^2)^4 = 81m^4n^8$
  \item $(m^{3a} \cdot m^{7a})^3 = (m^{10a})^3 = m^{30a}$
  \item $\big[(3x)^2 \cdot (2x^3)^2\big]^3 = (9x^2 \cdot 4x^6)^3 = (36x^8)^3 = 36^3x^{24} = 46\,656\,x^{24}$
  \item $\left[y^2 \cdot (3y^2)^2\right] \div 9y^4 = (y^2 \cdot 9y^4) \div 9y^4 = 9y^6 \div 9y^4 = y^2$
  \item $\left(\dfrac{a^{2x}}{a^3}\right)^3 = \left(a^{2x-3}\right)^3 = a^{6x-9}$
  \item $\left(\dfrac{k^{7t}}{k^{3t}}\right)^{10} = \left(k^{4t}\right)^{10} = k^{40t}$
  \item $\left(\dfrac{w^{5m}}{w^m}\right)^{-1} = \left(w^{4m}\right)^{-1} = w^{-4m} = \dfrac{1}{w^{4m}}$
  \item $\left(\dfrac{p^{10x}}{p^{6x}}\right)^{-3} = \left(p^{4x}\right)^{-3} = p^{-12x} = \dfrac{1}{p^{12x}}$
  \item $\left(\dfrac{a^{8m} \cdot a^{2m}}{a^{4m}}\right)^6 = \left(a^{6m}\right)^6 = a^{36m}$
  \item $\left(\dfrac{x^{4a} \cdot x^{5b}}{x^{2a} \cdot x^{3b}}\right)^6 = \left(x^{2a+2b}\right)^6 = x^{12a+12b}$
  \item $\left(\dfrac{n^{5x}}{n^{3x}} \cdot \dfrac{n^{7x}}{n^{3x}}\right)^6 = \left(n^{2x} \cdot n^{4x}\right)^6 = \left(n^{6x}\right)^6 = n^{36x}$
  \item En base $2$: $64 = 2^6$ y $128 = 2^7$, así que $64^{5x} = 2^{30x}$ y $128^x = 2^{7x}$. La división: $2^{30x} \div 2^{7x} = 2^{23x}$, y al elevar a $5$: $\left(2^{23x}\right)^5 = 2^{115x}$.
  \item En base $3$: $27 = 3^3$ y $9 = 3^2$, así que $27^p = 3^{3p}$ y $9^{3p} = 3^{6p}$. El producto: $3^{3p} \cdot 3^{6p} = 3^{9p}$, y al elevar a $6$: $\left(3^{9p}\right)^6 = 3^{54p}$.
\end{enumerate}

\vspace{10pt}
\begin{cuadrosolucion}
  \textbf{\textquestiondown Cómo te fue?} Ya dominas las potencias y hasta usaste letras como los algebristas: productos de potencias con exponentes literales ($a^{7x} \cdot a^{6x}$), exponentes negativos ($a^{-5}$), cocientes de potencias y hasta \textbf{cambio de base} ($64 = 2^6$, $27 = 3^3$). Eso es el corazón del álgebra. En la guía 4 daremos el gran salto: las \textbf{ecuaciones}, donde las letras se convierten en misterios que debemos resolver. \textquestiondown Allá nos vemos!
\end{cuadrosolucion}

\end{document}
